% !PRINTED 447
\weberpage{447}
Daraus ergibt sich nun wieder nach (7), daß die Kongruenz % [sic] Der Druck nennt (7); die Folgerung stützt sich auf die unmittelbar vorangehenden Kongruenzen (19)--(20).
\weq{21}{\xi^{p^\varepsilon} \equiv \xi \quad (\mathrm{mod}\ \mathfrak P)}
für jede ganze Zahl $\xi$ in $\mathfrak K$ erfüllt ist und daraus also der Satz:

\gesperrt{Eine Primzahl $p$ zerfällt im Körper $\mathfrak K$ in lauter voneinander verschiedene Primideale, deren Grad gleich dem Index von $p$ ist.}

Wir wollen noch beweisen, daß es unter den verschiedenen
in $p$ aufgehenden Primidealen $\mathfrak P$ immer eines gibt, für welches
von den beiden aus (16) folgenden Kongruenzen die Kongruenz
(17) besteht.

Dazu bemerken wir: Wenn wir die sämtlichen Zahlen des
Körpers $\mathfrak K$ in die konjugiert imaginären verwandeln, indem wir
$\sqrt{D}$ mit $-\sqrt{D}$ und jede Klasseninvariante $(k)$ mit der entgegengesetzten $(k^{-1})$ vertauschen, so gehen die sämtlichen Zahlen eines
Ideals $\mathfrak A$ in die Zahlen eines konjugiert imaginären Ideals $\mathfrak A'$
über, das auch mit $\mathfrak A$ identisch sein kann, und die Norm des
Ideals $\mathfrak A$ ist gleich der Norm des Ideals $\mathfrak A'$. Ist $\mathfrak A$ ein Primideal,
so ist auch $\mathfrak A'$ ein Primideal.

Ist nun die Kongruenz (17) nicht erfüllt, so muß nach (16)
\[
(lk)^p \equiv (k) \quad (\mathrm{mod}\ \mathfrak P)
\]
sein, und auch diese Kongruenz bleibt bestehen, wenn $k$ durch
eine andere Klasse $k'$ ersetzt wird. Setzen wir also $l^{-1}k^{-1}$ an
Stelle von $k$, so folgt:
\[
(k^{-1})^p \equiv (l^{-1}k^{-1}) \quad (\mathrm{mod}\ \mathfrak P).
\]

Da also $(k^{-1})^p - (l^{-1}k^{-1})$ eine Zahl in $\mathfrak P$ ist, so ist $(k)^p - (lk)$
in dem zu $\mathfrak P$ konjugierten Ideal $\mathfrak P'$ enthalten, und es ist
\[
(k)^p \equiv (lk) \quad (\mathrm{mod}\ \mathfrak P'),
\]
was aus (17) durch Vertauschung von $\mathfrak P$ mit $\mathfrak P'$ hervorgeht.

7. Aus diesen Sätzen ist die Irreducibilität der Klassengleichung auch nach Adjunktion von $\sqrt{D}$ eine einfache Folge.

Sind $k, k'$ zwei beliebige Klassen der Diskriminante $D$, so
können wir $k'$ nach § 118 in der Weise zusammensetzen:
\[
k' = k l l' l'' \ldots,
\]
daß durch die Klassen $l, l', l'', \ldots$ die Primzahlen $p, p', p'', \ldots$
darstellbar sind. Nach dem, was soeben bewiesen ist, können
wir also die Primteiler $\mathfrak P, \mathfrak P', \mathfrak P''$ dieser Primzahlen so wählen, daß
%CONT>>
